\chapter{1995:Paul J. Crutzen and Mario J. Molina and F. Sherwood Rowland}

\label{chap:chemistry1995}

The Nobel Prize in Chemistry for the year 1995 was awarded jointly to Paul J. Crutzen and Mario J. Molina and F. Sherwood Rowland.

\textbf{Motivation:}

for their work in atmospheric chemistry, particularly concerning the formation and decomposition of ozone

\section{Paul J. Crutzen}

\subsection*{Early Life and Education}

Paul J. Crutzen was born in 1933 in Amsterdam, the Netherlands.

\subsection*{Career and Major Research}

Crutzen earned his Ph.D. from the University of Stockholm in 1973. He became a director at the Max Planck Institute for Chemistry in Mainz, where he worked from 1980 to 2000.

\subsection*{Nobel Prize-winning Work}

The work for which Paul J. Crutzen was awarded the Nobel Prize in Chemistry in 1995 was described by the Nobel Committee as: "for their work in atmospheric chemistry, particularly concerning the formation and decomposition of ozone".

He showed that nitrogen oxides catalyze the destruction of stratospheric ozone, clarifying the chemical cycles that control the ozone layer.

\subsection*{Significance and Impact}

Establishing the role of nitrogen oxides in ozone chemistry was key to predicting how human activity could affect the ozone layer.

\subsection*{Later Career and Honors}

Crutzen continued his research at the Max Planck Institute for Chemistry in Mainz. He shared the 1995 Nobel Prize in Chemistry, and he died in 2021.

\subsection*{Legacy}

His work is fundamental to atmospheric chemistry, and he helped introduce the term Anthropocene for the current epoch of human influence on the Earth system.

\subsection*{Selected Publications}

"The Influence of Nitrogen Oxides on the Atmospheric Ozone Content" (Quarterly Journal of the Royal Meteorological Society, 1970)

\clearpage

\section{Mario J. Molina}

\subsection*{Early Life and Education}

Mario J. Molina was born in 1943 in Mexico City, Mexico.

\subsection*{Career and Major Research}

Molina studied chemical engineering at the Universidad Nacional Autonoma de Mexico (1965) and earned his Ph.D. from the University of California, Berkeley, in 1972. He was a postdoctoral researcher with F. Sherwood Rowland at the University of California, Irvine, and later held professorships at the Massachusetts Institute of Technology and the University of California, San Diego.

\subsection*{Nobel Prize-winning Work}

The work for which Mario J. Molina was awarded the Nobel Prize in Chemistry in 1995 was described by the Nobel Committee as: "for their work in atmospheric chemistry, particularly concerning the formation and decomposition of ozone".

He showed that chlorofluorocarbons are photolyzed in the stratosphere, releasing chlorine atoms that catalytically destroy ozone.

\subsection*{Significance and Impact}

The CFC-ozone hypothesis, developed with F. Sherwood Rowland, led to international regulation of chlorofluorocarbons under the Montreal Protocol.

\subsection*{Later Career and Honors}

Molina continued research at the University of California, San Diego, and later founded the Molina Center for Energy and the Environment in Mexico City. He shared the 1995 Nobel Prize in Chemistry, and he died in 2020.

\subsection*{Legacy}

His work linked atmospheric science directly to international policy and contributed to the recovery of the ozone layer.

\subsection*{Selected Publications}

"Stratospheric Sink for Chlorofluoromethanes: Chlorine Atom-Catalysed Destruction of Ozone" (Nature, 1974)

\clearpage

\section{F. Sherwood Rowland}

\subsection*{Early Life and Education}

F. Sherwood Rowland was born in 1927 in Delaware, Ohio, USA.

\subsection*{Career and Major Research}

Rowland studied at Ohio Wesleyan University (B.A., 1948) and earned his Ph.D. from the University of Chicago in 1952. He joined the University of California, Irvine, in 1964.

\subsection*{Nobel Prize-winning Work}

The work for which F. Sherwood Rowland was awarded the Nobel Prize in Chemistry in 1995 was described by the Nobel Committee as: "for their work in atmospheric chemistry, particularly concerning the formation and decomposition of ozone".

With Mario J. Molina, he established that chlorofluorocarbons accumulate in the stratosphere and decompose to release chlorine that destroys ozone.

\subsection*{Significance and Impact}

The prediction was confirmed by measurements of the Antarctic ozone hole, which led to the Montreal Protocol banning ozone-depleting substances.

\subsection*{Later Career and Honors}

Rowland remained at the University of California, Irvine, where he was Donald Bren Professor of Chemistry. He shared the 1995 Nobel Prize in Chemistry, and he died in 2012.

\subsection*{Legacy}

His atmospheric research demonstrated the global consequences of long-lived industrial gases.

\subsection*{Selected Publications}

"Stratospheric Sink for Chlorofluoromethanes: Chlorine Atom-Catalysed Destruction of Ozone" (Nature, 1974)

\clearpage

\section*{Chapter Summary}

The 1995 Nobel Prize in Chemistry recognized atmospheric chemistry and the formation and decomposition of ozone. Crutzen identified the role of nitrogen oxides, while Molina and Rowland showed that chlorofluorocarbons deplete stratospheric ozone.
